Full gate + dropout $-$ matched early & +0.0354 & 0.0206 & Overall pre-specified system comparison\\
Gate without dropout $-$ matched early & +0.0449 & 0.0341 & Separate stems plus dynamic weighting\\
Full gate + dropout $-$ gate without dropout & -0.0095 & 0.0258 & Dropout increment inside gated architecture\\
Early + dropout $-$ matched early & +0.0038 & 0.0322 & Dropout increment inside early fusion\\
Gate without dropout $-$ fixed equal stems & +0.0621 & 0.0188 & Dynamic weighting beyond equal stem fusion\\
Gate without dropout $-$ learned projection & +0.0404 & 0.0074 & Dynamic weighting beyond deterministic learned fusion\\
\bottomrule
\end{tabularx}
\end{table*}

\subsection{Image-Level Bootstrap Uncertainty}
An archived conditional uncertainty analysis resamples validation images 2,000 times with a fixed bootstrap seed. Each resample computes the project-local metric for fixed seed-0 and seed-2 checkpoints, averages those two checkpoints within each model group, and then calculates the reliability-minus-early difference. Images with no ground-truth boxes remain in the resampling frame; they contribute false positives when predictions are present and zero ground-truth boxes to denominators. The percentile intervals for AP50, AP75, and F1 are positive (\tabref{tab:bootstrap}). These intervals are conditional on two fixed checkpoints and this validation partition; they do not replace the three-seed training variability above and are not used for model selection.

\begin{table}[t]
\centering
\caption{Image-level bootstrap differences, reliability-aware $p=0.15$ minus matched early fusion. Values are descriptive percentile intervals from 2,000 image resamples.}
\label{tab:bootstrap}
\scriptsize
\begin{tabular}{@{}lcc@{}}
\toprule
Metric & Mean difference & 95\% CI\\
\midrule
AP50 & +0.0177 & [+0.0153, +0.0203]\\
AP75 & +0.0553 & [+0.0483, +0.0622]\\
F1 & +0.0187 & [+0.0150, +0.0226]\\
\bottomrule
\end{tabular}
\end{table}

\subsection{Synthetic Channel Removal}
To test controlled missing-input behavior, one channel group is deterministically set to zero at a time while weights, thresholds, and postprocessing remain fixed. \tabref{tab:removal} summarizes the archived seed-0/seed-2 checkpoint means. The reliability-aware system retains higher AP50 after removal of RGB, thermal, or event inputs. The largest separation follows thermal removal, where AP50 is 0.7030 for reliability-aware fusion and 0.3648 for early fusion. This deterministic intervention does not emulate measured physical sensor outages, temporal misalignment, or cross-device deployment failures.

\begin{table*}[t]
\centering
\caption{Synthetic channel removal on component-disjoint validation. Each value is the two-checkpoint mean. Delta is relative to the all-modal mean of the same model group. The intervention is deterministic zero-channel input, not a physical sensor-failure experiment.}
\label{tab:removal}
\scriptsize
\begin{tabular}{@{}llrrrrr@{}}
\toprule
Model group & Condition & Precision & Recall & F1 & AP50 & AP75\\
\midrule
Matched early & All modalities & 0.9188 & 0.8720 & 0.8945 & 0.9408 & 0.8208\\
Matched early & RGB removed & 0.8202 & 0.6531 & 0.7269 & 0.7379 & 0.5265\\
Matched early & Thermal removed & 0.8226 & 0.3003 & 0.4400 & 0.3648 & 0.2626\\
Matched early & Event removed & 0.9001 & 0.8731 & 0.8863 & 0.9339 & 0.8021\\
\midrule
Reliability $p=0.15$ & All modalities & 0.9273 & 0.8997 & 0.9133 & 0.9586 & 0.8760\\
Reliability $p=0.15$ & RGB removed & 0.8569 & 0.8229 & 0.8394 & 0.8985 & 0.7051\\
Reliability $p=0.15$ & Thermal removed & 0.8321 & 0.5885 & 0.6894 & 0.7030 & 0.3883\\
Reliability $p=0.15$ & Event removed & 0.9305 & 0.8970 & 0.9134 & 0.9582 & 0.8600\\
\bottomrule
\end{tabular}
\end{table*}

\begin{figure}[t]
\centering
\includegraphics[width=\columnwidth]{figures/fig4_channel_removal.pdf}
\caption{AP50 under deterministic synthetic channel removal. Reliability-aware $p=0.15$ retains higher AP50 under all four input conditions.}
\label{fig:removal}
\end{figure}

\subsection{Trained Single-Modality Baselines}
To quantify the standalone information in each sensor stream, RGB-only, thermal-only, and event-only detectors were trained with the same frozen component-disjoint train/validation manifests for seeds 0, 1, and 2. The supplied V76 records report thresholded precision, recall, and F1 together with AP50 and AP75. COCO AP@[0.50:0.95] and AR were not supplied for these nine runs and are therefore not reconstructed.

\begin{table*}[t]
\centering
\caption{Three-seed trained single-modality results on component-disjoint TriAir validation.}
\label{tab:single-modal-seeds}
\scriptsize
\begin{tabular}{@{}lrrrrrr@{}}
\toprule
Modality & Seed & Precision & Recall & F1 & AP50 & AP75\\
\midrule
RGB-only & 0 & 0.812 & 0.572 & 0.671 & 0.645 & 0.372\\
RGB-only & 1 & 0.825 & 0.590 & 0.687 & 0.662 & 0.389\\
RGB-only & 2 & 0.818 & 0.581 & 0.679 & 0.651 & 0.381\\
Thermal-only & 0 & 0.851 & 0.749 & 0.797 & 0.842 & 0.616\\
Thermal-only & 1 & 0.862 & 0.765 & 0.811 & 0.858 & 0.638\\
Thermal-only & 2 & 0.856 & 0.757 & 0.804 & 0.849 & 0.625\\
Event-only & 0 & 0.682 & 0.286 & 0.403 & 0.322 & 0.118\\
Event-only & 1 & 0.701 & 0.305 & 0.425 & 0.347 & 0.134\\
Event-only & 2 & 0.690 & 0.297 & 0.415 & 0.335 & 0.126\\
\bottomrule
\end{tabular}
\end{table*}

\begin{table*}[t]
\centering
\caption{Single-modality three-seed means and sample standard deviations.}
\label{tab:single-modal-summary}
\scriptsize
\setlength{\tabcolsep}{3pt}
\begin{tabular}{@{}lrrrrr@{}}
\toprule
Modality & Precision & Recall & F1 & AP50 & AP75\\
\midrule
RGB-only & $0.8183\pm0.0065$ & $0.5810\pm0.0090$ & $0.6790\pm0.0080$ & $0.6527\pm0.0086$ & $0.3807\pm0.0085$\\
Thermal-only & $\mathbf{0.8563\pm0.0055}$ & $\mathbf{0.7570\pm0.0080}$ & $\mathbf{0.8040\pm0.0070}$ & $\mathbf{0.8497\pm0.0080}$ & $\mathbf{0.6263\pm0.0111}$\\
Event-only & $0.6910\pm0.0095$ & $0.2960\pm0.0095$ & $0.4143\pm0.0110$ & $0.3347\pm0.0125$ & $0.1260\pm0.0080$\\
\bottomrule
\end{tabular}
\end{table*}

Thermal is the strongest standalone modality on every reported mean, while event-only is the weakest. The full reliability-aware system exceeds thermal-only by $0.1057\pm0.0133$ F1, $0.1037\pm0.0094$ AP50, and $0.2465\pm0.0253$ AP75 in seed-paired comparisons; all nine metric differences are positive across seeds. Matched early fusion also exceeds thermal-only by $0.0876\pm0.0101$ AP50 and $0.1827\pm0.0352$ AP75. These comparisons establish that the multimodal gain is not explained solely by selecting the strongest thermal stream, while the poor event-only result indicates that the event representation is most useful as complementary rather than standalone evidence under this protocol.

\subsection{Locked Internal Holdout}
After the six matched early and full reliability-aware checkpoints were fixed, the unchanged 837-image guard manifest was evaluated once without retraining, threshold tuning, or checkpoint reselection. The split inventory contains 2,287 boxes, while the standardized evaluator retains 1,264 ground-truth boxes after its label and geometry contract; both counts are reported to avoid conflating inventory and evaluator populations.

\tabref{tab:guard} shows that reliability-aware fusion improves mean AP50 from $0.9588\pm0.0058$ to $0.9674\pm0.0033$. The paired AP50 gain is $0.0086\pm0.0062$ and is positive for all three seeds. Mean F1 increases by $0.0069\pm0.0089$. AP75 improves slightly on average but is mixed by seed, including a negative seed-2 difference. The guard therefore provides useful locked internal holdout evidence, not an independent public test or a basis for post-hoc tuning.
